%\begin{center}
%\large \bf \runtitle
%\end{center}
%\vspace{1cm}
\chapter*{\runtitle}

\noindent \emph{Timbre morphing}, the generation of intermediate timbres
between known instruments, is a tool of interest for music production
and sound design, as it allows the sonic palette to be expanded beyond
traditional instruments. Common techniques operate in the latent space
of a single generative model trained on multiple instruments; while a
high-capacity model can represent several timbres with fidelity, in small
models this strategy compromises the fidelity of each individual timbre.

This work proposes an alternative approach: training variational
autoencoders ($\beta$-VAE) specialized per instrument from a common base
model, and generating hybrid timbres by interpolating the weights of those
models.

The system is trained on magnitude spectrograms obtained via STFT
from four instruments (piano, guitar, voice and bass) and is evaluated
using Fréchet Audio Distance (FAD), cosine similarity over MERT
\emph{embeddings}, latent-space analysis with UMAP and reconstruction
error.

\bigskip

\noindent\textbf{Keywords:} timbre morphing, variational autoencoder,
$\beta$-VAE, weight interpolation, audio synthesis, Fréchet Audio
Distance.
